% Source-aligned English for printed p.200.
\clearpage
\SourcePage{200}

\LayerTarget{tohoku:setup:equivariant-derived-functors}
We now restrict $A$ to the category ${\mathbf C}^{X(G)}$ of abelian $G$-sheaves on
$X$. Then $\Gamma_X^G$ is a left-exact additive functor from ${\mathbf C}^{X(G)}$ to
${\mathbf C}$, the category of abelian groups, and $f_*^G$ is a left-exact functor from
${\mathbf C}^{X(G)}$ to ${\mathbf C}^Y$, the category of abelian sheaves (without operators) on
$Y$. We set

\LayerTarget{tohoku:equation:5.2.3}
\[
H^n(X;G,A)=R^n\Gamma_X^G(A)
\tag{5.2.3}
\]

\LayerTarget{tohoku:equation:5.2.4}
\[
\mathbf H^n(G,A)=R^nf_*^G(A)
\tag{5.2.4}
\]
\LayerTarget{tohoku:result:equivariant-cohomology-universality}
Thus, for an abelian $G$-sheaf, the $H^n(X;G,A)$ are abelian groups and the
$\mathbf H^n(G,A)$ are sheaves on $Y$. Both form universal cohomological
functors in $A$, reducing for $n=0$ to $\Gamma_X^G(A)$ and $A^G$,
respectively. If $G$ acts trivially on $X$, we recover the notions introduced
in 4.4. --- As in Proposition 3.1.3, one also sees readily that if $U$ is an
open subset of $Y$ and $A$ is an \textit{injective} abelian $G$-sheaf on $X$, then its
restriction to $f^{-1}(U)$ is an injective abelian $G$-sheaf on that space.
Hence, as in 4.2:

\LayerTarget{tohoku:proposition:5.2.1}
\textsc{Proposition 5.2.1.} \textit{Let $A$ be an abelian
$G$-sheaf on $X$. Then, for every open subset $U\subset Y$, $\mathbf H^p(G,A)\vert U=\mathbf H^p(G,A\vert f^{-1}(U))$.}

\LayerTarget{tohoku:remark:sheaf-group-cohomology-notation}
This is why $X$ is omitted from the notation $\mathbf H^p(G,A)$. A more
substantial reduction in the calculation of the $\mathbf H^p(G,A)$ is given
by the following result.

\LayerTarget{tohoku:corollary:5.2.1}
\textsc{Corollary.} \textit{Suppose that $f^{-1}(U)$ is the union
of pairwise disjoint open subsets $g\mathbin{.}V$ ($g\in G/G_0$), where
$V$ is open in $X$ and $G_0$ is a subgroup of $G$ such that
$g_0\mathbin{.}V=V$ for $g_0\in G_0$. Then $\mathbf H^p(G,A)\vert U=\mathbf H^p(G_0,A\vert V)$,
under the natural identification $U=V/G_0$.}

\LayerTarget{tohoku:proof:corollary-5.2.1}
By Proposition 5.2.1, we may assume $U=Y$, whence $f^{-1}(U)=X$. The category
of $G$-sheaves on $X$ is then isomorphic to the category of $G_0$-sheaves on
$V$, and this isomorphism is compatible with the functors $f_*^G$ and
$f_*^{G_0}$. The stated formula follows at once. --- A more explicit
calculation will be given below (cf. formula (5.2.11) and Theorem 5.3.1).

\LayerTarget{tohoku:construction:two-compositions-of-gamma-x-g}
The formulas (5.2.2) express $\Gamma_X^G$ as a composite functor in two
different ways. In both cases Theorem 2.4.1 applies, by virtue of the
corollary to Proposition 5.1.3. We therefore obtain:

\LayerTarget{tohoku:theorem:5.2.1}
\textsc{Theorem 5.2.1.} \textit{On the category ${\mathbf C}^{X(G)}$ of
abelian $G$-sheaves on $X$, there are two cohomological spectral functors
converging to the graded functor $(H^n(X;G,A))$, whose initial terms are,
respectively,}

\LayerTarget{tohoku:equation:5.2.5}
\[
\tag{5.2.5}
\left\{
\begin{aligned}
\mathrm I_2^{p,q}(A)&=H^p(Y,\mathbf H^q(G,A))\\
\mathrm {II}_2^{p,q}(A)&=H^p(G,H^q(X,A)).
\end{aligned}
\right.
\]

\LayerTarget{tohoku:proof:theorem-5.2.1:p1}
To establish the explicit form of $\mathrm {II}_2^{p,q}(A)$, it remains only
to show that the derived functors of $\Gamma_X$, viewed as a functor from
${\mathbf C}^{X(G)}$ to ${\mathbf C}^G$, the category of $G$-modules, are indeed the
$H^q(X,A)$. This follows immediately

\TohokuContinuation{201}
